\documentclass[12pt,a4paper]{article}

% === Packages ===
\usepackage[UTF8]{ctex}
\usepackage{geometry}
\geometry{margin=1in}
\usepackage{amsmath,amssymb}
\usepackage{listings}
\usepackage{xcolor}
\usepackage{hyperref}
\usepackage{graphicx}
\usepackage{booktabs}
\usepackage{enumitem}
\usepackage{fancyhdr}
\usepackage{tocloft}

% === Code listing style ===
\definecolor{codebg}{rgb}{0.95,0.95,0.95}
\definecolor{codeframe}{rgb}{0.8,0.8,0.8}
\lstset{
    basicstyle=\ttfamily\small,
    backgroundcolor=\color{codebg},
    frame=single,
    rulecolor=\color{codeframe},
    breaklines=true,
    numbers=none,
    tabsize=4,
    language=C++,
    keywordstyle=\color{blue},
    commentstyle=\color{gray},
    stringstyle=\color{red}
}

% === Title ===
\title{\textbf{Contributing to OpenCV\\
        Fixing Signed Left-Shift Undefined Behavior\\
        in Drawing Code}}
\author{李语尚\\\small Student ID: 12412308}
\date{CS219: Advanced Programming\\Project 4\\2026 年 5 月 30 日}

\begin{document}
\maketitle
\thispagestyle{empty}

\newpage
\tableofcontents
\newpage

% ======================================================================
\section{背景与任务选择}
% ======================================================================

\subsection{Project 4 简介}

本次 project 要求向开源项目 OpenCV 提交一个有意义的贡献。OpenCV 是世界上最广泛使用的计算机视觉库之一，主要使用 C++ 编写，被学术界和工业界广泛采用。

在 Project 4 的三个等级中（Level A：初学者、Level B：中级、Level C：高级），我选择了 \textbf{Level C — 解决真实的 OpenCV Issue}。具体而言，我修复了 OpenCV issue \href{https://github.com/opencv/opencv/issues/28940}{\#28940}：绘图模块中负坐标导致的有符号整数左移 undefined behavior。

\subsection{任务概述}

问题出在 \texttt{imgproc} 模块的绘图实现中。当用户调用 \texttt{cv::drawContours()} 并传入负的 \texttt{offset} 参数时，轮廓点可能被移动到图像边界外形成负坐标。负坐标本身是合法输入（表示图形部分位于图像外，需要裁剪），但原始代码在 fixed-point 坐标转换时对负的有符号整数执行左移，触发了 C++ undefined behavior。

\begin{itemize}[leftmargin=*]
    \item \textbf{Issue}: \href{https://github.com/opencv/opencv/issues/28940}{\texttt{\#28940}}
    \item \textbf{涉及文件}: \texttt{modules/imgproc/src/drawing.cpp}
    \item \textbf{修复方式}: 将 signed left shift 替换为乘法缩放
    \item \textbf{影响范围}: \texttt{ThickLine()}, \texttt{FillConvexPoly()}, \texttt{CollectPolyEdges()}
\end{itemize}

% ======================================================================
\section{问题分析}
% ======================================================================

\subsection{触发场景}

一个现实场景是：用户从原图中裁剪出 ROI，然后希望把原图坐标系中的轮廓绘制到 ROI 图像中。例如：

\begin{lstlisting}
cv::Mat roi(80, 80, CV_8UC3, cv::Scalar::all(0));

std::vector<std::vector<cv::Point>> contours = {{
    cv::Point(100, 100), cv::Point(160, 100),
    cv::Point(160, 160), cv::Point(100, 160)
}};

cv::Point offset(-120, -120);

cv::drawContours(roi, contours, -1, cv::Scalar(0, 255, 0),
                 2, cv::LINE_8, cv::noArray(), INT_MAX, offset);
\end{lstlisting}

第一个轮廓点 $(100, 100)$ 加上 offset $(-120, -120)$ 后变成了 $(-20, -20)$。这是合法的输入——OpenCV 应该裁剪后绘制可见部分。

\subsection{根因分析}

\texttt{drawing.cpp} 中定义了 fixed-point 坐标缩放常量：

\begin{lstlisting}
enum { XY_SHIFT = 16, XY_ONE = 1 << XY_SHIFT };
\end{lstlisting}

绘图模块内部使用 16.16 定点数格式，因此会将普通坐标放大 $2^{16}$ 倍。原始代码在 \texttt{ThickLine()} 中使用了 signed 左移：

\begin{lstlisting}
p0.x <<= XY_SHIFT - shift;
p0.y <<= XY_SHIFT - shift;
\end{lstlisting}

当 \texttt{p0.x == -20} 时，等价于执行 \texttt{-20 << 16}。根据 C++ 标准，对负的 signed integer 进行左移是 \textbf{undefined behavior}。普通 Release 构建下，x86 平台使用补码表示，\emph{可能}看起来可以运行；但 C++ 标准不保证这种行为，且 UBSan 会在运行时检测并终止程序。

此外，粗线绘制会将线段扩展成四边形后调用 \texttt{FillConvexPoly()} 填充——第一次只修复 \texttt{ThickLine()} 后，同一输入会继续触发 \texttt{FillConvexPoly()} 中的同类问题。因此需要修复同一绘图路径上的所有 fixed-point 转换。

\subsection{受影响的函数}

以下三个函数均可能接收由负 offset 产生的负坐标：

\begin{enumerate}[leftmargin=*]
    \item \texttt{ThickLine()} —— 粗线绘制，约 line 1669–1672
    \item \texttt{FillConvexPoly()} —— 凸多边形填充，约 line 1128–1150
    \item \texttt{CollectPolyEdges()} —— 多边形边收集，约 line 1274–1321
\end{enumerate}

% ======================================================================
\section{修复方案}
% ======================================================================

\subsection{设计原则}

\begin{enumerate}[leftmargin=*]
    \item \textbf{不禁止负坐标} —— 负坐标是合法输入，表示图形部分位于图像外。
    \item \textbf{不 clamp 到 0} —— 这会导致图形形状变形，绘制结果不正确。
    \item \textbf{保持原 fixed-point 缩放语义} —— $x \ll 16 \approx x \times 65536$。
    \item \textbf{避免 signed negative left shift UB} —— 使用乘法替代。
\end{enumerate}

\subsection{新增 helper 函数}

\begin{lstlisting}
// Use multiplication to avoid undefined behavior from
// left-shifting negative coordinates.
static inline int64
ScaleToFixedPoint( int64 v, int shift )
{
    CV_Assert( 0 <= shift && shift <= XY_SHIFT );
    const int64 factor = static_cast<int64>(1) << (XY_SHIFT - shift);
    CV_DbgAssert( v >= std::numeric_limits<int64>::min() / factor &&
                  v <= std::numeric_limits<int64>::max() / factor );
    return v * factor;
}
\end{lstlisting}

关键设计要点：
\begin{itemize}[leftmargin=*]
    \item \texttt{factor} 是 $2^{\text{XY\_SHIFT} - \text{shift}}$，始终为非负整数。左移非负数在 C++ 中是定义良好的。
    \item 返回值使用乘法：\texttt{v * factor}。对于非负数，结果与左移一致；对于负数，乘法是 C++ 定义良好的行为。
    \item \texttt{CV\_DbgAssert} 在 Debug 模式下检测溢出，保证安全。
\end{itemize}

\subsection{修改范围（Before / After）}

\subsubsection{ThickLine()}

\textbf{Before:}
\begin{lstlisting}
p0.x <<= XY_SHIFT - shift;
p0.y <<= XY_SHIFT - shift;
p1.x <<= XY_SHIFT - shift;
p1.y <<= XY_SHIFT - shift;
\end{lstlisting}

\textbf{After:}
\begin{lstlisting}
p0.x = ScaleToFixedPoint(p0.x, shift);
p0.y = ScaleToFixedPoint(p0.y, shift);
p1.x = ScaleToFixedPoint(p1.x, shift);
p1.y = ScaleToFixedPoint(p1.y, shift);
\end{lstlisting}

\subsubsection{FillConvexPoly()}

\textbf{Before:}
\begin{lstlisting}
p0.x <<= XY_SHIFT - shift;
p0.y <<= XY_SHIFT - shift;
// ...
p.x <<= XY_SHIFT - shift;
p.y <<= XY_SHIFT - shift;
// ...
xs <<= XY_SHIFT - shift;
xe <<= XY_SHIFT - shift;
\end{lstlisting}

\textbf{After:}
\begin{lstlisting}
p0.x = ScaleToFixedPoint(p0.x, shift);
p0.y = ScaleToFixedPoint(p0.y, shift);
// ...
p.x = ScaleToFixedPoint(p.x, shift);
p.y = ScaleToFixedPoint(p.y, shift);
// ...
xs = ScaleToFixedPoint(xs, shift);
xe = ScaleToFixedPoint(xe, shift);
\end{lstlisting}

\subsubsection{CollectPolyEdges()}

\textbf{Before:}
\begin{lstlisting}
pt0.x = (pt0.x + offset.x) << (XY_SHIFT - shift);
// ...
pt1.x = (pt1.x + offset.x) << (XY_SHIFT - shift);
// ...
pt0c.x = (int64)(t0.x) << XY_SHIFT;
pt1c.x = (int64)(t1.x) << XY_SHIFT;
// ...
t0.y = pt0.y << XY_SHIFT;
t1.y = pt1.y << XY_SHIFT;
\end{lstlisting}

\textbf{After:}
\begin{lstlisting}
pt0.x = ScaleToFixedPoint(pt0.x + offset.x, shift);
// ...
pt1.x = ScaleToFixedPoint(pt1.x + offset.x, shift);
// ...
pt0c.x = ScaleToFixedPoint(t0.x, 0);
pt1c.x = ScaleToFixedPoint(t1.x, 0);
// ...
t0.y = ScaleToFixedPoint(pt0.y, 0);
t1.y = ScaleToFixedPoint(pt1.y, 0);
\end{lstlisting}

\subsection{方案比较}

我也考虑了其他修复方案：

\textbf{方案 A: 使用乘法缩放（已采用）}
\begin{lstlisting}
return v * (static_cast<int64>(1) << (XY_SHIFT - shift));
\end{lstlisting}
优点：语义清晰，对负数定义良好，易于解释，性能影响可忽略。

\textbf{方案 B: 转为 unsigned 后左移}
\begin{lstlisting}
return static_cast<int64>(static_cast<uint64_t>(v) << shift);
\end{lstlisting}
优点：避免 signed left shift UB，与补码行为接近。\\
缺点：unsigned 转 signed 在极端值上容易引发维护者讨论，可读性不如乘法。

\textbf{方案 C: 负坐标 clamp 到 0}\\
不推荐。负坐标是合法输入，clamp 会改变几何形状导致绘制不正确。

综上，采用 \textbf{方案 A}。

\subsection{回归测试}

在 \texttt{modules/imgproc/test/test\_drawing.cpp} 中新增回归测试：

\begin{lstlisting}
TEST(Drawing, drawContours_negative_offset_regression_28940)
{
    Mat roi(80, 80, CV_8UC1, Scalar::all(0));
    vector<vector<Point>> contours(1);
    contours[0].push_back(Point(100, 100));
    contours[0].push_back(Point(160, 100));
    contours[0].push_back(Point(160, 160));
    contours[0].push_back(Point(100, 160));

    drawContours(roi, contours, -1, Scalar(255), 2,
                 LINE_8, noArray(), INT_MAX, Point(-120, -120));

    // 轮廓移动到 (-20,-20) 到 (40,40)，只有右下部分可见
    EXPECT_NE(0, roi.at<uchar>(20, 40));
    EXPECT_NE(0, roi.at<uchar>(40, 20));
    EXPECT_NE(0, roi.at<uchar>(40, 40));
    EXPECT_EQ(0, countNonZero(roi(Rect(60, 60, 20, 20))));
}
\end{lstlisting}

该测试覆盖了：
\begin{enumerate}[leftmargin=*]
    \item \textbf{可见区域}：$(20, 40)$, $(40, 20)$, $(40, 40)$ 应该有像素被绘制——验证裁剪后的轮廓仍然正确落在 ROI 内。
    \item \textbf{越界区域}：ROI 右下角 $(60, 60)$ 到 $(79, 79)$ 应该全为 0——验证没有越界写入。
\end{enumerate}

% ======================================================================
\section{验证与测试}
% ======================================================================

\subsection{UBSan 构建与复现}

为了在修复前确认 UB 的存在并在修复后验证它已被消除，我在 WSL 中构建了一个 UBSan 版本的 OpenCV，仅包含 \texttt{core} 和 \texttt{imgproc} 模块：

\begin{lstlisting}[language=bash]
cmake -S . -B build_ubsan_28940 \
  -DCMAKE_BUILD_TYPE=Debug \
  -DBUILD_LIST=core,imgproc \
  -DCMAKE_CXX_FLAGS="-fsanitize=undefined \
    -fno-sanitize-recover=undefined \
    -fno-omit-frame-pointer"

cmake --build build_ubsan_28940 --target opencv_imgproc -j8
\end{lstlisting}

使用复现程序 \texttt{test\_28940.cpp} 调用 \texttt{drawContours()}，分别加载原始库和修复后库。

\subsection{修复前}

\begin{lstlisting}[language=bash]
$ LD_PRELOAD=./build_ubsan_28940/lib/libopencv_imgproc.so \
  ./test_28940

Calling drawContours with offset [-120, -120].
  The first point becomes [-20, -20]

drawing.cpp:1662:10:
  runtime error: left shift of negative value -20

#0 ThickLine drawing.cpp:1662
#1 cv::drawContours drawing.cpp:2571
#2 main test_28940.cpp:41
\end{lstlisting}

UBSan 明确检测到 \texttt{left shift of negative value}，证实了 undefined behavior 的存在。

\subsection{修复后}

\begin{lstlisting}[language=bash]
$ LD_PRELOAD=./build_ubsan_fixed/lib/libopencv_imgproc.so \
  ./test_28940

Calling drawContours with offset [-120, -120].
  The first point becomes [-20, -20]

drawContours finished, pixel sum = [0, 61710, 0, 0]
\end{lstlisting}

修复后：
\begin{enumerate}[leftmargin=*]
    \item \texttt{drawContours()} 正常完成绘制。
    \item UBSan 不再报告错误。
    \item 输出图像非空——像素和 $= [0, 61710, 0, 0]$，验证图像确实被绘制了。
\end{enumerate}

% ======================================================================
\section{性能评估}
% ======================================================================

\subsection{测试方法}

为了确认修复不会明显影响正常使用性能，我分别使用原始 Release 库和修复后 Release 库运行 benchmark。测试配置如下：

\begin{itemize}[leftmargin=*]
    \item 构建类型：Release
    \item 不启用 UBSan
    \item 模块：\texttt{core, imgproc}
    \item 每组重复 5 次，记录 best 和 avg
    \item 输出 checksum 保持一致以确认绘制结果相同
\end{itemize}

\subsection{正常正坐标数据（offset = (0, 0)）}

\begin{table}[h]
\centering
\caption{正常坐标下 best time 对比}
\begin{tabular}{rrrr}
\toprule
\textbf{迭代次数} & \textbf{原始 best (ms)} & \textbf{修复 best (ms)} & \textbf{差异} \\
\midrule
100   & 32.56 & 31.84 & -2.21\% \\
1000  & 327.38 & 321.65 & -1.75\% \\
10000 & 3228.03 & 3252.81 & +0.77\% \\
\bottomrule
\end{tabular}
\end{table}

\begin{table}[h]
\centering
\caption{正常坐标下 avg time 对比}
\begin{tabular}{rrrr}
\toprule
\textbf{迭代次数} & \textbf{原始 avg (ms)} & \textbf{修复 avg (ms)} & \textbf{差异} \\
\midrule
100   & 32.89 & 32.29 & -1.84\% \\
1000  & 328.55 & 326.96 & -0.48\% \\
10000 & 3274.25 & 3268.45 & -0.18\% \\
\bottomrule
\end{tabular}
\end{table}

输出 checksum 一致：$1.19493 \times 10^7$。

\subsection{负 offset 数据（offset = (-120, -120)）}

\begin{table}[h]
\centering
\caption{负 offset 下 best time 对比}
\begin{tabular}{rrrr}
\toprule
\textbf{迭代次数} & \textbf{原始 best (ms)} & \textbf{修复 best (ms)} & \textbf{差异} \\
\midrule
100   & 30.33 & 30.79 & +1.51\% \\
1000  & 304.72 & 305.12 & +0.13\% \\
10000 & 2997.30 & 3033.59 & +1.21\% \\
\bottomrule
\end{tabular}
\end{table}

\begin{table}[h]
\centering
\caption{负 offset 下 avg time 对比}
\begin{tabular}{rrrr}
\toprule
\textbf{迭代次数} & \textbf{原始 avg (ms)} & \textbf{修复 avg (ms)} & \textbf{差异} \\
\midrule
100   & 30.67 & 31.34 & +2.19\% \\
1000  & 307.94 & 309.08 & +0.37\% \\
10000 & 3023.27 & 3056.36 & +1.09\% \\
\bottomrule
\end{tabular}
\end{table}

输出 checksum 一致：$1.07391 \times 10^7$。

\subsection{性能结论}

修复后的性能与原始版本基本一致，差异大多在 0\%–2\% 范围内，属于正常运行噪声。原因有三：

\begin{enumerate}[leftmargin=*]
    \item 修改只发生在少量几何点的 fixed-point 坐标转换处。
    \item 绘图函数主要耗时在裁剪、扫描线填充和像素写入。
    \item \texttt{v * (1 << k)} 是乘以 2 的幂，Release 编译器通常会优化为等价的低成本指令（LEA 或位移 + 加法）。
\end{enumerate}

因此，本修复基本不影响性能。

% ======================================================================
\section{修改文件总结}
% ======================================================================

\begin{table}[h]
\centering
\caption{修改文件清单}
\begin{tabular}{lll}
\toprule
\textbf{文件} & \textbf{修改性质} & \textbf{行数变化} \\
\midrule
\texttt{modules/imgproc/src/drawing.cpp} & 新增 \texttt{ScaleToFixedPoint} 函数 + 替换 17 处 left shift & +28, -18 \\
\texttt{modules/imgproc/test/test\_drawing.cpp} & 新增回归测试 & +20 \\
\bottomrule
\end{tabular}
\end{table}

% ======================================================================
\section{困难与反思}
% ======================================================================

\subsection{遇到的困难}

\begin{enumerate}[leftmargin=*]
    \item \textbf{UBSan 一次报错不全}：最初只修复 \texttt{ThickLine()} 后，以为已经解决了问题。但重新运行测试后发现粗线绘制内部调用的 \texttt{FillConvexPoly()} 也有同样的 signed left shift。这说明在大型代码库中，同一类 bug 经常存在于多个位置，需要系统性地排查整个调用链。
    
    \item \textbf{Replace 模式的保守性}：OpenCV 是一个非常成熟的项目，任何对核心函数的修改都必须谨慎。我最初的修改方案在 \texttt{CV\_DbgAssert} 中使用了除法来检测溢出，需要确保除法不会引入新的性能问题（好在 Debug 模式不追求极致性能）。
    
    \item \textbf{回归测试的边界选择}：写一个既简洁又能覆盖关键路径的测试并不容易。需要确保测试点落在可见区域和不可见区域的边界上，且不受具体绘制实现细节（如像素级偏移）的影响。
\end{enumerate}

\subsection{AI 工具使用}

在本项目中，我使用了 \textbf{CodeWhale（DeepSeek V4）} 作为 AI 辅助编程工具。具体用途包括：

\begin{itemize}[leftmargin=*]
    \item \textbf{代码库分析}：阅读和搜索 OpenCV 的 \texttt{drawing.cpp}（约 2900 行），定位所有 signed left shift 的位置。CodeWhale 的并行搜索能力在这个阶段非常高效——它能同时 grep 多个符号并交叉验证。
    \item \textbf{修复方案推演}：讨论了「乘法」vs「unsigned cast」两种方案，选择更易维护和解释的乘法方案。
    \item \textbf{Benchmark 代码编写}：生成了 \texttt{bench\_28940.cpp} 和 \texttt{bench\_28940\_negative.cpp}，用于系统的 before/after 性能对比。
    \item \textbf{测试代码编写}：帮助设计了回归测试的验证逻辑和 assert 条件。
    \item \textbf{报告撰写}：协助整理和结构化本报告的初稿。
\end{itemize}

与之前 project 中 AI 主要是“写代码”不同，这次 AI 更像是“结对编程的资深工程师”——帮助理解大型项目的代码结构、设计方案、验证思路。真正的决策（选择什么方案、如何测试、提交什么内容）仍然由我自己做出。

\subsection{收获与体会}

这次 project 最大的收获不是学会了某个具体技术，而是体验了一次完整的 \textbf{真实开源贡献流程}：

\begin{enumerate}[leftmargin=*]
    \item \textbf{从零构建一个大型项目}：OpenCV 的 CMake 构建系统非常复杂，需要理解模块依赖、编译选项、sanitizer 集成等。与之前的 homework 项目不同，这不是从头写一个小程序，而是在一个千万行级别的代码库中做手术式的修改。
    
    \item \textbf{阅读和理解他人代码}：\texttt{drawing.cpp} 中包含了 Bresenham 画线、多边形扫描线填充、抗锯齿等图形学算法。虽然不需要完全理解每一步，但需要理解数据流——坐标如何从外部传入、经过哪些转换、最终如何写入图像。
    
    \item \textbf{C++ undefined behavior 的可怕之处}：Release 模式下 \texttt{-20 << 16} 在 x86 上恰好“工作正常”，但这是被标准明确定义为 UB 的操作。编译器完全有权生成不正确的代码。UBSan 是这类问题的照妖镜——没有它，这类 bug 可能会潜伏数年。
    
    \item \textbf{开源社区的工作方式}：从 issue 报告到 PR 提交，每一步都需要清晰的沟通、可复现的证据、对审核者的尊重。修复代码是一部分，写好 commit message 和 PR 描述是同样重要的一部分。
    
    \item \textbf{性能验证的重要性}：在像 OpenCV 这样的性能敏感项目中，即使是修复 bug 也需要提供性能数据证明没有退化。这迫使我建立了系统的 benchmark 流程。
\end{enumerate}

Project-based learning 确实是成为真正程序员的最佳途径。这次 project 让我从“写代码的人”向“软件工程师”迈进了一步。

% ======================================================================
\section{总结}
% ======================================================================

本次 project 修复了 OpenCV \texttt{imgproc} 绘图模块在负坐标场景下的 C++ undefined behavior。修复将 signed left shift 替换为乘法缩放，在保持语义不变的前提下消除了 UB。

\textbf{修复前}：
\begin{center}
\texttt{drawContours(..., offset=(-120,-120))} \\
$\rightarrow$ \texttt{ThickLine()} $\rightarrow$ \texttt{p0.x <<= 16} \\
$\rightarrow$ \texttt{left shift of negative value} \\
$\rightarrow$ \textbf{UBSan 报错，程序终止}
\end{center}

\textbf{修复后}：
\begin{center}
\texttt{drawContours(..., offset=(-120,-120))} \\
$\rightarrow$ fixed-point 坐标转换使用乘法缩放 \\
$\rightarrow$ \textbf{无 UBSan 错误} \\
$\rightarrow$ \textbf{图像正常绘制，checksum 一致}
\end{center}

该修复保持绘制结果与原始版本一致，性能影响在 0\%–2\% 范围内可以忽略，适合作为一个独立的、真实的、可测试的 OpenCV bug fix PR。

\subsection*{GitHub 仓库}

\begin{itemize}[leftmargin=*]
    \item \textbf{Fork}: \url{https://github.com/lys0204/opencv}
    \item \textbf{Branch}: \texttt{project4\_qdB3ar}
    \item \textbf{Commit}: \texttt{434bb82a8878} --- \texttt{imgproc: avoid signed left-shift UB in drawing code}
\end{itemize}

% ======================================================================
\begin{thebibliography}{9}

\bibitem{issue28940}
OpenCV Issue \#28940: \textit{[UB] Left Shift of Negative Value in cv::ThickLine When Drawing Contours with Negative Offset}.\\
\url{https://github.com/opencv/opencv/issues/28940}

\bibitem{contributing}
OpenCV Contribution Guide.\\
\url{https://github.com/opencv/opencv/wiki/How\_to\_contribute}

\bibitem{cppstandard}
ISO/IEC 14882:2020 --- Programming Language C++.\\
Section 7.6.7: Shift operators. Left shift of a negative signed integer is undefined behavior.

\bibitem{ubsan}
LLVM Documentation: UndefinedBehaviorSanitizer.\\
\url{https://clang.llvm.org/docs/UndefinedBehaviorSanitizer.html}

\bibitem{opencvrepo}
OpenCV GitHub Repository.\\
\url{https://github.com/opencv/opencv}

\bibitem{myfork}
My OpenCV Fork.\\
\url{https://github.com/lys0204/opencv}

\end{thebibliography}

\end{document}
